% !TeX root = ../main.tex
\chapter{Dualità e applicazione trasposta}
\label{cap:dualita}
\fonti{Lezioni 19, 20 e 21}

Un'applicazione lineare a valori nel campo associa uno scalare a ogni vettore. Lo spazio di queste applicazioni è il duale: consente di descrivere coordinate ed equazioni con un linguaggio comune, e di interpretare la trasposizione delle matrici come un'operazione sulle applicazioni.

Gli spazi sono sullo stesso campo $K$; nei risultati sulle dimensioni si dichiarerà l'ipotesi di dimensione finita. Ricordiamo che $\operatorname{Hom}_K(V,W)$ è lo spazio delle applicazioni lineari da $V$ a $W$, con somma e prodotto per scalari definiti punto per punto, e che $[v]_{\mathcal B}$ è la colonna delle coordinate nella base ordinata $\mathcal B$.
% !TeX root = ../main.tex
\section{Spazio duale e base duale}
\begin{definizione}
Il \emph{duale} di uno spazio vettoriale $V$ su $K$ è
\[
 V^*=\operatorname{Hom}_K(V,K)=\{\varphi:V\to K:\varphi\text{ è lineare}\}.
\]
I suoi elementi sono detti \emph{funzionali lineari} o \emph{forme lineari}. Lo zero di $V^*$ è la funzione che vale $0_K$ su ogni vettore; le operazioni sono $(\varphi+\psi)(v)=\varphi(v)+\psi(v)$ e $(\lambda\varphi)(v)=\lambda\varphi(v)$.
\end{definizione}
\begin{definizione}
Sia $\mathcal B=(v_1,\ldots,v_n)$ una base di $V$. Il funzionale $v_i^*\in V^*$ legge la coordinata $i$:
\[
 v_i^*\!\left(\sum_{j=1}^n\lambda_jv_j\right)=\lambda_i.
\]
In particolare $v_i^*(v_j)=\delta_{ij}$, dove il \emph{delta di Kronecker} è lo scalare
\[
 \delta_{ij}=\begin{cases}1,&i=j,\\0,&i\ne j.\end{cases}
\]
La famiglia ordinata $\mathcal B^*=(v_1^*,\ldots,v_n^*)$ si chiama \emph{base duale} di $\mathcal B$.
\end{definizione}
La funzione $v_i^*$ è ben definita per l'unicità delle coordinate ed è lineare perché la coordinata $i$ di $av+bw$ è la combinazione con coefficienti $a,b$ delle coordinate $i$ di $v,w$. Un elemento $v_i^*$ è una funzione, non un vettore di $V$ ottenuto modificando $v_i$.

\Needspace{12\baselineskip}
\begin{teorema}
La base duale è una base di $V^*$. Se $\dim V=n$, allora $\dim V^*=n$, e per ogni $\varphi\in V^*$ vale
\[
 \varphi=\sum_{i=1}^n\varphi(v_i)v_i^*.
\]
\end{teorema}
\begin{proof}
\textbf{Indipendenza.} Se $\sum_i\alpha_iv_i^*=0$ come funzionale, valutando sul vettore $v_j$ si ottiene
\[
 0=\sum_i\alpha_iv_i^*(v_j)=\sum_i\alpha_i\delta_{ij}=\alpha_j.
\]
Questo vale per ogni $j$, quindi tutti i coefficienti sono nulli.

\textbf{Generazione.} Per $\varphi\in V^*$, i coefficienti dell'eventuale espressione nella base duale devono essere $\varphi(v_i)$, come mostra la stessa valutazione. Poniamo dunque $\psi=\sum_i\varphi(v_i)v_i^*$. Le applicazioni $\psi$ e $\varphi$ coincidono su ogni $v_j$, quindi coincidono su tutto $V$: una funzione lineare è determinata dalle immagini di una base. La famiglia duale genera ed è indipendente, e perciò ha dimensione $n$.
\end{proof}

\begin{esempio}[Vettori come colonne e funzionali come righe]
Per $V=\R^n$ nella base canonica $\mathcal B=(e_1,\ldots,e_n)$, un vettore è una colonna $x=(x_1,\ldots,x_n)^T$ e un funzionale ha la forma
\[
 \varphi(x)=a_1x_1+\cdots+a_nx_n
 =\begin{pmatrix}a_1&\cdots&a_n\end{pmatrix}
 \begin{pmatrix}x_1\\\vdots\\x_n\end{pmatrix}.
\]
La matrice di $\varphi:\R^n\to\R$ è una riga di tipo $1\times n$. Invece la colonna delle sue coordinate \emph{nella base duale} è $[\varphi]_{\mathcal B^*}=(a_1,\ldots,a_n)^T$: riga della matrice e colonna delle coordinate contengono gli stessi scalari ma hanno ruoli diversi. Per esempio, in $\R^2$, il funzionale $\varphi(x_1,x_2)=2x_1-3x_2$ applicato a $(4,1)^T$ restituisce lo scalare $5$.
\end{esempio}

\section{Dualità canonica e spazio biduale}
\subsection{Valutare un funzionale su un vettore}
\begin{definizione}
La \emph{dualità canonica} fra $V$ e $V^*$ è la funzione di valutazione
\[
 V\times V^*\longrightarrow K,\qquad (v,\varphi)\longmapsto
 \langle v,\varphi\rangle=\varphi(v).
\]
È \emph{bilineare}, cioè lineare in ciascun argomento quando l'altro è fissato.
\end{definizione}
Infatti la linearità di $\varphi$ e le operazioni punto per punto in $V^*$ danno
\begin{align*}
 \langle av+bw,\varphi\rangle&=a\langle v,\varphi\rangle+b\langle w,\varphi\rangle,\\
 \langle v,a\varphi+b\psi\rangle&=a\langle v,\varphi\rangle+b\langle v,\psi\rangle.
\end{align*}
Qui i due argomenti appartengono a spazi diversi: la scrittura non presuppone un prodotto scalare su $V$. In una base e nella sua duale si ha $\langle v,\varphi\rangle=[\varphi]_{\mathcal B^*}^{\,T}[v]_{\mathcal B}$.

\subsection{L'isomorfismo canonico con il biduale}
\begin{definizione}
Il \emph{biduale} è $V^{**}=(V^*)^*=\operatorname{Hom}_K(V^*,K)$. Per ogni $v\in V$ definiamo
\[
 \alpha_v:V^*\longrightarrow K,\qquad \alpha_v(\varphi)=\varphi(v),
\]
e la funzione $\alpha:V\to V^{**}$ data da $\alpha(v)=\alpha_v$. Un elemento di $V^{**}$ riceve dunque un funzionale come argomento e restituisce uno scalare.
\end{definizione}
\begin{teorema}
Se $V$ ha dimensione finita, $\alpha:V\to V^{**}$ è un isomorfismo canonico, ossia definito senza scegliere una base.
\end{teorema}
\begin{proof}
\textbf{Buona definizione e linearità.} Per $v$ fissato, la linearità nel secondo argomento della dualità dà
\[
 \alpha_v(a\varphi+b\psi)=a\alpha_v(\varphi)+b\alpha_v(\psi),
\]
quindi $\alpha_v\in V^{**}$. Inoltre, per ogni $\varphi\in V^*$,
\[
 \alpha_{av+bw}(\varphi)=\varphi(av+bw)
 =a\alpha_v(\varphi)+b\alpha_w(\varphi).
\]
Si tratta dell'uguaglianza di due funzionali su $V^*$: pertanto $\alpha$ è lineare.

\textbf{Iniettività.} Scelta una base $\mathcal B=(v_1,\ldots,v_n)$, scriviamo $v=\sum_i\lambda_iv_i$. Se $\alpha(v)=0$, valutando sui funzionali della base duale otteniamo
\[
 0=\alpha_v(v_i^*)=v_i^*(v)=\lambda_i\qquad(i=1,\ldots,n).
\]
Quindi $v=0$ e il nucleo di $\alpha$ è nullo.

\textbf{Suriettività.} Applicando due volte il risultato sulla dimensione del duale si trova $\dim V^{**}=\dim V^*=\dim V=n$. Una funzione lineare iniettiva fra spazi della stessa dimensione finita è suriettiva, dunque $\alpha$ è un isomorfismo. La base è servita nella prova, ma la formula $\alpha_v(\varphi)=\varphi(v)$ non ne usa alcuna: per questo l'isomorfismo è canonico.
\end{proof}
\begin{schema}
\begin{tikzpicture}[font=\small]
 \node (v) at (0,0) {$v\in V$};
 \node (a) at (5,0) {$\alpha_v\in V^{**}$};
 \node (ph) at (5,-1.3) {$\varphi\in V^*$};
 \node (s) at (9,-1.3) {$\varphi(v)\in K$};
 \draw[-{Stealth},coloreschema] (v)--node[above] {$\alpha$} (a);
 \draw[-{Stealth},coloreschema] (ph)--node[below] {$\alpha_v$} (s);
\end{tikzpicture}

$\alpha$ trasforma un vettore nell'operazione «valuta un funzionale su quel vettore».
\end{schema}

\begin{osservazione}
In dimensione finita anche $V$ e $V^*$ sono isomorfi, perché hanno la stessa dimensione, ma la sola struttura di spazio vettoriale non specifica un isomorfismo privilegiato fra loro. Una base ne permette la scelta; su uno spazio reale di dimensione finita, un prodotto scalare $g$ fornisce invece la mappa $v\mapsto\varphi_v$, con $\varphi_v(w)=g(v,w)$.\footnote{Qui un prodotto scalare reale è una funzione bilineare simmetrica $g:V\times V\to\R$ tale che $g(v,v)>0$ per $v\ne0$. La mappa indicata è lineare e iniettiva, perché $g(v,w)=0$ per ogni $w$ implica $g(v,v)=0$ e quindi $v=0$. In dimensione finita è allora un isomorfismo, per l'uguaglianza delle dimensioni.} Senza aver scelto $g$, questa formula non determina una funzione. L'identificazione con il biduale, al contrario, è data direttamente dalla valutazione.

In dimensione infinita le conclusioni sulle dimensioni non si trasferiscono. Per il duale algebrico, con le usuali ipotesi di scelta, si ha $\dim V<\dim V^*<\dim V^{**}$. Le dimensioni sono qui intese come cardinalità di basi; questo risultato generale è citato senza dimostrazione; non viene usato nelle prove in dimensione finita.
\end{osservazione}

% !TeX root = ../main.tex
\Needspace{20\baselineskip}
\section{Annullatori di sottospazi}
\begin{definizione}
Per un sottospazio $U\subseteq V$ si definisce il suo \emph{annullatore}
\[
 U^\perp=\{\varphi\in V^*: \varphi(u)=0\text{ per ogni }u\in U\}\subseteq V^*.
\]
Per un sottospazio $L\subseteq V^*$ si definisce invece
\[
 L^\perp=\{v\in V:\varphi(v)=0\text{ per ogni }\varphi\in L\}\subseteq V.
\]
Questi spazi sono detti anche \emph{ortogonali} rispetto alla dualità canonica. Il simbolo $\perp$ non presuppone un prodotto scalare: per capire in quale spazio si trova l'annullatore bisogna guardare se si parte da vettori o da funzionali.
\end{definizione}
\begin{proposizione}
\begin{enumerate}[label=\arabic*.]
\item $U^\perp$ è un sottospazio di $V^*$.
\item $L^\perp$ è un sottospazio di $V$.
\end{enumerate}
\end{proposizione}
\begin{proof}
\begin{enumerate}[label=\arabic*.]
\item \passomath{U^\perp} Il funzionale nullo annulla $U$. Per $\varphi,\psi\in U^\perp$ e $a,b\in K$, ogni $u\in U$ soddisfa $(a\varphi+b\psi)(u)=a\varphi(u)+b\psi(u)=0$. Quindi $a\varphi+b\psi\in U^\perp$.
\item \passomath{L^\perp} Il vettore nullo è annullato da tutti gli elementi di $L$. Se $v,w\in L^\perp$, allora $\varphi(av+bw)=a\varphi(v)+b\varphi(w)=0$ per ogni $\varphi\in L$. Pertanto $av+bw\in L^\perp$.
\end{enumerate}
\end{proof}

\begin{teorema}[Dimensione dell'annullatore]
\label{teo:dimensione-annullatore}
Se $\dim V=n<\infty$ e $\dim U=r$, allora $\dim U^\perp=n-r$.
\end{teorema}
\begin{proof}
Scegliamo una base $(u_1,\ldots,u_r)$ di $U$ e completiamola a una base
\[
 \mathcal B=(u_1,\ldots,u_r,v_{r+1},\ldots,v_n)
\]
di $V$. La base duale è $\mathcal B^*=(u_1^*,\ldots,u_r^*,v_{r+1}^*,\ldots,v_n^*)$. Ogni $\varphi\in V^*$ si scrive in modo unico come
\[
 \varphi=\sum_{i=1}^r a_iu_i^*+\sum_{j=r+1}^n b_jv_j^*.
\]
Annullare tutto $U$ equivale ad annullare ciascun vettore della sua base, cioè a imporre $\varphi(u_i)=a_i=0$ per $i=1,\ldots,r$. Ne segue
\[
 U^\perp=\operatorname{span}\{v_{r+1}^*,\ldots,v_n^*\}.
\]
Questi $n-r$ funzionali sono indipendenti perché fanno parte di una base; costituiscono quindi una base di $U^\perp$. Nei casi estremi $U=\{0\}$ e $U=V$ si ottengono rispettivamente $U^\perp=V^*$ e $U^\perp=\{0\}$.
\end{proof}

\begin{proposizione}[Doppio annullatore]
Se $V$ ha dimensione finita, allora $(U^\perp)^\perp=U$, dove il secondo annullatore è inteso come sottospazio di $V$.
\end{proposizione}
\begin{proof}
\passomath{\supseteq} Per definizione ogni funzionale in $U^\perp$ si annulla su ogni vettore di $U$, quindi $U\subseteq(U^\perp)^\perp$.

\passomath{\subseteq} Se $v\notin U$, una base $(u_1,\ldots,u_r)$ di $U$, seguita da $v$, è indipendente e può essere completata a una base di $V$. Il funzionale della base duale corrispondente a $v$ vale $1$ su $v$ e $0$ su tutti gli $u_i$. Appartiene dunque a $U^\perp$ ma non annulla $v$: pertanto $v\notin(U^\perp)^\perp$. Nessun vettore esterno a $U$ appartiene al doppio annullatore.
\end{proof}
Questa verifica esplicita il modo in cui un funzionale può distinguere un vettore esterno da un sottospazio. Se il secondo annullatore viene invece costruito nel biduale, l'uguaglianza va letta come identificazione mediante l'isomorfismo canonico $\alpha$.

\Needspace{13\baselineskip}
\begin{esempio}
In $V=\R^3$, sia $U=\operatorname{span}\{(1,1,0)\}$. Un funzionale generico è $\varphi(x_1,x_2,x_3)=ax_1+bx_2+cx_3$; appartiene a $U^\perp$ se e solo se $a+b=0$. Quindi
\[
 U^\perp=\operatorname{span}\{\varphi_1,\varphi_2\},\qquad
 \varphi_1(x)=x_1-x_2,\quad\varphi_2(x)=x_3.
\]
Si vede sia $\dim U^\perp=2=3-1$, sia il recupero del sottospazio mediante le equazioni:
\[
 (U^\perp)^\perp=\{x:x_1-x_2=0,\ x_3=0\}
 =\{(t,t,0):t\in\R\}=U.
\]
Le equazioni cartesiane di un sottospazio possono così essere interpretate come funzionali che lo annullano.
\end{esempio}

% !TeX root = ../main.tex
\section{La trasposta di un'applicazione lineare}
\begin{definizione}
Per $f:V\to W$ lineare, la sua \emph{trasposta} è
\[
 f^*:W^*\longrightarrow V^*,\qquad f^*(\psi)=\psi\circ f.
\]
Il verso si inverte: un funzionale su $W$ diventa un funzionale su $V$ precomponendolo con $f$. Qui l'asterisco indica la trasposta algebrica, senza coniugazione complessa e senza scegliere un prodotto scalare.
\end{definizione}
\begin{schema}
\begin{tikzpicture}[font=\small]
 \node (v) at (0,0) {$V$};\node (w) at (4,0) {$W$};\node (k) at (2,-1.6) {$K$};
 \draw[-{Stealth},coloreschema] (v)--node[above] {$f$} (w);
 \draw[-{Stealth}] (w)--node[below right] {$\psi$} (k);
 \draw[-{Stealth},coloreschema] (v)--node[below left] {$\psi\circ f=f^*(\psi)$} (k);
 \node (ws) at (9,0) {$W^*$};\node (vs) at (6.5,0) {$V^*$};
 \draw[-{Stealth},coloreschema] (ws)--node[above] {$f^*$} (vs);
 \node[align=center] at (7.75,-1.2) {un funzionale su $W$\\diventa un funzionale su $V$};
\end{tikzpicture}
\end{schema}
La funzione $f^*(\psi)$ è lineare perché è una composizione di funzioni lineari. Anche $f^*$ è lineare: per $a,b\in K$ e $\psi,\eta\in W^*$,
\[
 f^*(a\psi+b\eta)=(a\psi+b\eta)\circ f=a(\psi\circ f)+b(\eta\circ f).
\]
L'uguaglianza fra funzionali si verifica valutando i due membri su ciascun $v\in V$.

\begin{proposizione}
La trasposta è caratterizzata dall'identità
\[
 \langle v,f^*(\psi)\rangle=\langle f(v),\psi\rangle
 \qquad(v\in V,\ \psi\in W^*).
\]
\end{proposizione}
\begin{proof}
Per definizione di valutazione e di composizione,
\[
 \langle v,f^*(\psi)\rangle=f^*(\psi)(v)=(\psi\circ f)(v)
 =\psi(f(v))=\langle f(v),\psi\rangle.
\]
Viceversa, se un funzionale ha questi valori per ogni $v$, coincide con $\psi\circ f$: la formula determina univocamente $f^*(\psi)$.
\end{proof}

\Needspace{20\baselineskip}
\subsection{Nucleo e immagine della trasposta}
\begin{teorema}
\begin{enumerate}[label=\arabic*.]
\item Per ogni applicazione lineare $f:V\to W$,
\[
 \ker f^*=(\operatorname{Im}f)^\perp\subseteq W^*.
\]
\item Se $V,W$ hanno dimensione finita, vale inoltre
\[
 \operatorname{Im}f^*=(\ker f)^\perp\subseteq V^*,\qquad
 \operatorname{rg}f^*=\operatorname{rg}f.
\]
\end{enumerate}
\end{teorema}
\begin{proof}
\begin{enumerate}[label=\arabic*.]
\item \textbf{Nucleo.} Per $\psi\in W^*$, la condizione $\psi\in\ker f^*$ significa che $f^*(\psi)$ è il funzionale nullo su $V$, cioè
\[
 \psi(f(v))=0\quad\text{per ogni }v\in V.
\]
Questo equivale a dire che $\psi$ annulla tutti i vettori di $\operatorname{Im}f$, ossia che $\psi\in(\operatorname{Im}f)^\perp$. Le due implicazioni danno l'uguaglianza.
\item \textbf{Inclusione per l'immagine.} Se $\varphi=f^*(\psi)$ e $u\in\ker f$, allora
\[
 \varphi(u)=\psi(f(u))=\psi(0_W)=0.
\]
Dunque $\operatorname{Im}f^*\subseteq(\ker f)^\perp$.

\textbf{Uguaglianza delle dimensioni.} Supponiamo ora $\dim V=n$, $\dim W=m$ e poniamo $r=\dim\operatorname{Im}f$. Dal risultato sul nucleo e dalla dimensione dell'annullatore,
\[
 \dim\ker f^*=m-r.
\]
Applicando il teorema della dimensione a $f^*:W^*\to V^*$ si trova
\[
 \dim\operatorname{Im}f^*=\dim W^*-\dim\ker f^*=m-(m-r)=r.
\]
D'altra parte $\dim\ker f=n-r$, quindi
\[
 \dim(\ker f)^\perp=n-\dim\ker f=n-(n-r)=r.
\]
L'inclusione già dimostrata è fra sottospazi della stessa dimensione finita e perciò è un'uguaglianza. Il conteggio dà anche $\operatorname{rg}f^*=\operatorname{rg}f$.
\end{enumerate}
\end{proof}

\Needspace{37\baselineskip}
\section{La matrice della trasposta nelle basi duali}
Siano $\mathcal B=(v_1,\ldots,v_n)$ una base di $V$ e $\mathcal C=(w_1,\ldots,w_m)$ una base di $W$, con basi duali $\mathcal B^*$ e $\mathcal C^*$. Ricordiamo che la colonna $j$ di $M_{\mathcal B}^{\mathcal C}(f)$ contiene le coordinate di $f(v_j)$ in $\mathcal C$.
\begin{teorema}
Se $A=M_{\mathcal B}^{\mathcal C}(f)\in M(m,n,K)$, allora
\[
 M_{\mathcal C^*}^{\mathcal B^*}(f^*)=A^T\in M(n,m,K).
\]
\end{teorema}
\begin{proof}
Scriviamo le immagini dei vettori delle due basi:
\[
 f(v_j)=\sum_{i=1}^m a_{ij}w_i,\qquad
 f^*(w_i^*)=\sum_{j=1}^n b_{ji}v_j^*.
\]
Per leggere il coefficiente $b_{hi}$ valutiamo il secondo funzionale in $v_h$. Usando la dualità e la definizione di trasposta otteniamo
\[
 b_{hi}=f^*(w_i^*)(v_h)=w_i^*(f(v_h))
 =w_i^*\!\left(\sum_{\ell=1}^m a_{\ell h}w_\ell\right)=a_{ih}.
\]
Quindi l'entrata $(h,i)$ della matrice di $f^*$ è l'entrata $(i,h)$ di $A$, esattamente la definizione della matrice trasposta.
\end{proof}
La precisazione «nelle basi duali» è essenziale: scegliendo basi arbitrarie di $V^*,W^*$ intervengono anche matrici di cambiamento di base.

\Needspace{18\baselineskip}
\begin{esempio}
Sia $f:\R^2\to\R^3$ definita da
\[
 f(x_1,x_2)=(x_1+2x_2,\ 3x_2,\ -x_1+x_2),\qquad
 A=\begin{pmatrix}1&2\\0&3\\-1&1\end{pmatrix}
\]
nelle basi canoniche $\mathcal B=(e_1,e_2)$ del dominio e $\mathcal C=(e_1,e_2,e_3)$ del codominio. Se $\psi(y_1,y_2,y_3)=ay_1+by_2+cy_3$, allora
\[
 f^*(\psi)(x_1,x_2)=(a-c)x_1+(2a+3b+c)x_2.
\]
Nelle basi duali ciò diventa
\[
 [f^*(\psi)]_{\mathcal B^*}
 =\begin{pmatrix}a-c\\2a+3b+c\end{pmatrix}
 =\underbrace{\begin{pmatrix}1&0&-1\\2&3&1\end{pmatrix}}_{A^T}
 \begin{pmatrix}a\\b\\c\end{pmatrix}.
\]
Il dominio di $f^*$ ha dimensione $3$ e il codominio dimensione $2$, nel verso opposto a $f$. Anche le dimensioni della matrice si sono scambiate.
\end{esempio}

